\chapter{Fundamentação e trabalhos relacionados}\label{cap:relacionados}

\section{Representações neurais de cenas e superfícies}

Gerar uma imagem de uma cena a partir de uma câmera nova exige estimar o que existe ao longo de cada raio que parte dessa câmera. O NeRF de \citet{mildenhall2020nerf} representa essa informação por uma rede neural. Ao consultar uma posição no espaço, a rede fornece uma densidade volumétrica, isto é, um valor usado para estimar quanto aquela região contribui para bloquear o raio. A consulta também fornece uma cor que depende da posição e da direção de observação. Os valores de vários pontos ao longo do raio são combinados por renderização volumétrica para formar cada pixel da imagem.

O NeRF original foi proposto para sintetizar novas vistas. Sua densidade pode ser convertida em uma superfície, mas a formulação não representa diretamente a distância até essa superfície. NeuS, de \citet{wang2021neus}, e VolSDF, de \citet{yariv2021volsdf}, incorporam uma função de distância assinada à representação implícita. O sinal indica de que lado da superfície está um ponto e o módulo representa sua distância à superfície. Essa formulação favorece a extração de malhas, embora não garanta, por si só, que a geometria reconstruída coincida com a geometria real.

Na robótica, outro problema é atualizar uma representação tridimensional enquanto a câmera se desloca. O iMAP, de \citet{sucar2021imap}, e o NICE-SLAM, de \citet{zhu2022niceslam}, usam representações implícitas para estimar simultaneamente o mapa e a posição da câmera. Esses sistemas aproximam campos neurais de tarefas robóticas, mas uma cena para o Gazebo ainda precisa definir separadamente os modelos que serão visualizados, movidos e usados no cálculo de colisões.

\section{Estrutura de objetos e reconstrução panóptica}

Para consultar ou editar um objeto sem alterar o restante da cena, a representação precisa distinguir cada objeto do fundo e dos demais elementos. O Panoptic Neural Fields, proposto por \citet{kundu2022panoptic}, representa o fundo e cada objeto por redes separadas. Cada objeto ocupa uma caixa tridimensional orientada, e sua rede fornece cor e densidade em coordenadas locais. Essa organização permite remover, adicionar ou clonar objetos e alterar a posição ou a orientação de suas caixas antes de renderizar uma nova imagem. A edição ocorre na representação neural; o método não tem como saída principal malhas de colisão ou um mundo de simulador.

Assim como o Panoptic Neural Fields, outros trabalhos procuram manter informações de objetos de maneira coerente entre diferentes imagens. Eles tratam, contudo, etapas distintas do problema. O Panoptic Lifting, de \citet{siddiqui2023panoptic}, recebe imagens, poses e segmentações panópticas bidimensionais e resolve as inconsistências dos identificadores entre as vistas para formar uma representação panóptica tridimensional. O Contrastive Lift, de \citet{bhalgat2023contrastive}, combina características aprendidas em diferentes ritmos para consolidar a separação das instâncias. O PanoRecon, de \citet{wu2024panorecon}, reconstrói a partir de vídeo monocular em tempo real e mantém as instâncias ao longo da sequência. Esses métodos produzem representações nas quais os objetos podem ser identificados, mas não organizam diretamente os objetos como modelos e colisões de uma cena no Gazebo.

A comparação geométrica desta dissertação requer ainda um método que não dependa de um campo neural. O COLMAP reconstrói a geometria comparando imagens obtidas de diferentes posições. Seu componente de reconstrução estereoscópica multivista, ou MVS, procura correspondências entre as imagens para estimar profundidade e orientação local da superfície; depois, combina os mapas de profundidade em uma nuvem de pontos \citep{schonberger2016mvs}. Nesta dissertação, a nuvem do COLMAP e a malha extraída do DM-NeRF são comparadas à geometria de referência do DM-SR pelas mesmas distâncias e pelos mesmos F-scores. A aparência das imagens, a separação dos objetos e o funcionamento no Gazebo são avaliados separadamente.

O DM-NeRF, proposto por \citet{wang2023dmnerf}, é o método usado para representar a cena e identificar suas instâncias. Ele acrescenta ao NeRF um campo de objetos, aprendido com rótulos bidimensionais, e permite consultar um código de instância junto à cor e à densidade. O artigo original também manipula objetos no campo neural por consultas inversas antes da renderização. A partir dessa representação, esta dissertação investiga outra sequência: extrai uma malha, associa os IDs à sua topologia, separa os objetos, trata a superfície de apoio e gera uma cena para o Gazebo Classic. A contribuição concentra-se na implementação e na avaliação dessas etapas posteriores.

\section{Geração de ambientes para simulação}

Gerar um ambiente para simulação exige definir a geometria dos objetos e, conforme a aplicação, suas partes móveis e relações. Os trabalhos a seguir partem de imagens ou vídeos, mas diferem quanto ao uso de modelos previamente disponíveis e quanto ao tipo de cena produzido.

O URDFormer, de \citet{chen2024urdformer}, parte de uma imagem e produz uma hierarquia em \textit{Unified Robot Description Format} (URDF), formato que descreve elos e juntas. O método identifica partes, relações e articulações e usa ativos tridimensionais disponíveis previamente. Sua saída favorece a estrutura cinemática; ela não procura reproduzir, por reconstrução multivista, cada superfície observada.

O DRAWER, de \citet{xia2025drawer}, recebe vídeo e combina geometria implícita, aparência e articulações para produzir ambientes interativos. O GRS, de \citet{zook2025grs}, parte de observações RGB-D, segmenta a cena e recupera ativos para compor tarefas de simulação. Nos dois casos, o objetivo inclui comportamento ou estrutura funcional. A cadeia desta dissertação mantém como geometria os dados extraídos do DM-NeRF e mede o que ocorre ao convertê-los em modelos independentes; ela não substitui os objetos observados por ativos de um catálogo.

Em trabalho anterior, \citet{lemos2025digital} exploraram a descrição e reconstrução de ambientes com segmentação panóptica, incluindo separação de componentes, tratamento do piso e geração de uma cena editável no Gazebo. A presente dissertação amplia esse desenvolvimento ao fixar o DM-NeRF como método principal e submeter a cadeia a um protocolo quantitativo que distingue síntese, geometria, estrutura e funcionamento. Os resultados da publicação anterior não são combinados com as medições atuais.

\section{Comparação por características}

A Tabela~\ref{tab:comparacao_metodos} permite comparar os trabalhos por três perguntas: quais dados entram no método, que representação ele produz e se essa representação já corresponde a uma cena para simulação. A indicação de que o simulador não é a saída principal significa que o trabalho resolve síntese, reconstrução ou compreensão de cena, mas não entrega diretamente modelos com geometria de colisão. As linhas sobre URDFormer, DRAWER e GRS mostram uma estratégia diferente, baseada em hierarquias, articulações ou ativos. A última linha descreve a cadeia avaliada nesta dissertação.

\begin{table}[H]
\centering
\small
\renewcommand{\arraystretch}{1.12}
\caption{Comparação qualitativa entre trabalhos relacionados e a cadeia desta dissertação.}
\label{tab:comparacao_metodos}
\begin{tabularx}{\textwidth}{>{\raggedright\arraybackslash}p{2.55cm}>{\raggedright\arraybackslash}p{3.05cm}>{\raggedright\arraybackslash}X}
\toprule
\textbf{Trabalho} & \textbf{Entrada} & \textbf{Representação, estrutura e relação com simulação} \\
\midrule
Panoptic Neural Fields & imagens e predições 2D & campo neural separado em objetos e fundo, com edição antes da renderização; simulador não é a saída principal \\
Panoptic Lifting & RGB, poses e máscaras panópticas & campo neural panóptico 3D; não gera cena de simulação \\
PanoRecon & vídeo monocular & reconstrução online com instâncias; simulador não é o objetivo principal \\
DM-NeRF & RGB, poses e mapas 2D de instância & campo neural com IDs de objeto e manipulação antes da renderização; não gera cena diretamente \\
COLMAP & RGB e poses fixas & nuvem densa por MVS, sem instâncias ou cena de simulação \\
URDFormer & imagem RGB & hierarquia URDF com ativos predefinidos para simulação articulada \\
DRAWER & vídeo & ambiente interativo com geometria e articulações \\
GRS & RGB-D & segmentação, recuperação de ativos e tarefas de simulação \\
Cadeia avaliada nesta dissertação & RGB, poses e mapas 2D de instância & malha DM-NeRF, IDs, objetos separados e cena SDFormat no Gazebo Classic \\
\bottomrule
\end{tabularx}
\end{table}

A tabela não estabelece uma classificação de desempenho, pois as entradas e as saídas não são equivalentes. Panoptic Neural Fields, Panoptic Lifting, PanoRecon e DM-NeRF tratam a representação ou a identificação dos objetos. URDFormer, DRAWER e GRS se aproximam da geração de conteúdo funcional, com articulações ou ativos. A cadeia avaliada aqui ocupa a passagem entre esses grupos: parte dos IDs e da geometria do DM-NeRF, materializa objetos independentes e verifica sua utilização no Gazebo.
